%
% ---------- header -----------------------------------------------------------
%
% project       kaneton
%
% license       kaneton
%
% file          /home/mycure/kaneton/view/book/assignments/assignments - Practicals.tex
%
% created       julien quintard   [mon may 14 19:56:45 2007]
% updated       julien quintard   [thu feb 17 21:25:38 2011]
%

%
% ---------- setup ------------------------------------------------------------
%

%
% path
%

\def\path{../..}

%
% template
%

%
% ---------- header -----------------------------------------------------------
%
% project       kaneton
%
% license       kaneton
%
% file          /home/mycure/kaneton/view/template/book.tex
%
% created       julien quintard   [wed may 16 18:16:50 2007]
% updated       julien quintard   [sun jun  3 12:00:37 2007]
%

%
% class
%

\documentclass[10pt,a4wide]{book}

%
% packages
%

\usepackage[english]{babel}
\usepackage[T1]{fontenc}
\usepackage{a4wide}
\usepackage{fancyheadings}
\usepackage{multicol}
\usepackage{indentfirst}
\usepackage{graphicx}
\usepackage{color}
\usepackage{xcolor}
\usepackage{verbatim}

\usepackage{aeguill}

\usepackage[Lenny]{\path/package/fncychap}

\pagestyle{fancy}

%
% -
%

\renewcommand{\-}{\vspace{\parskip}}

%
% toc
%

\newcommand{\toc}
  {
    \tableofcontents
    \setlength{\footrulewidth}{0.3pt}
    \setlength{\parindent}{0.3cm}
    \setlength{\parskip}{2ex plus 0.5ex minus 0.2ex}
  }

%
% logos
%

\newcommand{\logos}
  {
    \begin{center}
      \includegraphics[scale=0.8]{\path/logo/kaneton.pdf}
    \end{center}
  }

%
% colors
%

\definecolor{functioncolor}{rgb}{0.40,0.00,0.00}
\definecolor{commandcolor}{rgb}{0.00,0.00,0.40}
\definecolor{verbatimcolor}{rgb}{0.00,0.40,0.00}
\definecolor{noticecolor}{rgb}{0.87,0.84,0.02}

%
% function
%

\newcommand\function[3]{
  \begin{tabular}{p{0.2cm}p{13.8cm}}
  & {\color{functioncolor}\textbf{#1}}#2
  \end{tabular}

  \begin{tabular}{p{1cm}p{13cm}}
  & #3
  \end{tabular}}

%
% align
%

\newcommand\align[1]{
  \\ & \hspace{#1}}

%
% argument
%

\newcommand\argument[1]{\textit{#1}}

%
% command
%

\newcommand\command[2]{
  \begin{tabular}{p{0.2cm}p{13.8cm}}
  & {\color{commandcolor}\textbf{#1}}
  \end{tabular}

  \begin{tabular}{p{1cm}p{13cm}}
  & #2
  \end{tabular}}

%
% notice
%

\newcommand\notice[1]{
  {\color{noticecolor}\textbf{Notice}}

  \begin{tabular}{p{0.2cm}p{13.8cm}}
  & #1
  \end{tabular}}

%
% subsubsubsection
%

\newcommand\subsubsubsection[1]{\textbf{#1}}

%
% example
%

\newcommand\example[1]{
  \textit{Example:}

  \begin{tabular}{p{0.2cm}p{13.8cm}}
  & \textit{#1}
  \end{tabular}}

%
% warning XXX
%

\renewcommand{\familydefault}{\sfdefault}

%
% verbatim stuff
%

\makeatletter

\renewcommand{\verbatim@font}
  {\ttfamily\footnotesize\selectfont}

\def\verbatim@processline{
  \hskip15ex{\color{verbatimcolor}\the\verbatim@line}\par
}

\makeatother

%
% header
%

\rhead{}
\rfoot{\scriptsize{The kaneton microkernel project}}

\date{\scriptsize{\today}}


%
% header
%

\lhead{\scriptsize{The kaneton microkernel :: Debug assignments}}

%
% title
%

\title{The kaneton microkernel :: Debug assignments \\
       \version
       \logo}

%
% document
%

\begin{document}

%
% title
%

\maketitle

%
% --------- text --------------------------------------------------------------
%

This document contains everything necessary for students to undertake the
kaneton educational project.

\-

All the kaneton documents are available on
  the official website\footnote{\location{http://kaneton.opaak.org}}.

%
% table of contents
%

\tableofcontents

%
% indentation
%

\indentation{}

%
% chapters
%

%
% ---------- header -----------------------------------------------------------
%
% project       kaneton
%
% license       kaneton
%
% file          /home/mycure/kaneton/view/book/assignments/k0.tex
%
% created       julien quintard   [fri nov 28 05:25:37 2008]
% updated       julien quintard   [sun mar  6 18:58:11 2011]
%


\chapter{Introduction}
\label{chapter:Introduction}

This practical session has two main goals.

First, the students should, by the end of the session, be able to manipulate their kernel with GDB.

The second advantage of this session is that students have an opportunity to start working on their own tests suite.

Everyone has to know that this session is the first of its kind since kaneton has been taught at EPITA. Thus, students must take the benefit of it. Moreover, they all are invited to report their comments on the subject in order to improve it for the next years.

\chapter{kaneton, Qemu and GDB}
\label{chapter:kaneton, Qemu and GDB}

So as to test their kaneton, students mostly use Qemu.

\section{Qemu}
The best way to understand what Qemu is capable of is to comment what is displayed on Qemu website :

QEMU is a generic and open source machine emulator and virtualizer.

\textit{When used as a machine emulator, Qemu can run OSes and programs made for one machine (e.g. an ARM board) on a different machine (e.g. your own PC). By using dynamic translation, it achieves very good performance.}

In our case, students use this function of Qemu. As an emulator, Qemu emulates hardware by including main of the hardware devices in a virtual environment. Then when the code is being executed, Qemu checks that opcode is compatible with the host's CPU. If not, Qemu converts it. That technique allow the execution on the host of binaries compiled for different architectures than the current one.

Obviously, the main default of emulation is that it is slow. For each instruction, the emulator has to check the compatibility and take the time to translate the incompatible instructions. On the other hand, the emulation is purely software. As a consequence, it is much more easy to use and much more flexible. The absence of hardware assistance is the main difference between emulation and virtualization.


\textit{When used as a virtualizer, QEMU achieves near native performances by executing the guest code directly on the host CPU. QEMU supports virtualization when executing under the Xen hypervisor or using the KVM kernel module in Linux. When using KVM, QEMU can virtualize x86, server and embedded PowerPC, and S390 guests.}

Contrary to emulation, virtualization consists in hardware ressources sharing with an hypervisor. KVM is a module for linux kernel that allows ressources sharing (hardware virtualization). This module is mandatory if one wants to use Qemu as a virtualizer. In addition to allow the execution of several OS in the same time, the couple KVM/Qemu improve the execution speed. The execution is here only twice slower than native solution, compared to emulation which is ten times slower.

\section{GDB}
GDB is a debugger. As so, the one who use it can define breakpoint, execute instructions step by step or modify on the fly some values like variables.

Naturally, a few flags must be added to gcc for the compilation. Those flags are -g -ggdb. The flag -g produces general debugging information (symbol for example) while -ggdb produces specific debugging information. Cf chapter "Options for Debugging Your Program or GCC" of GCC manual.

I won't develop the description of GDB more deeper. I consider that the reader has at least met this tool once. Nevertheless, it might be helpful to have a summary of the main commands :

\begin{verbatim}
(gdb) list : List the next lines of code.
(gdb) b main : Put a breakpoint on the beginning of the function main.
(gdb) c : Continue the execution to the next breakpoint.
(gdb) s : Next step, follow the function calls
(gdb) n : Next instruction, doesn't follow the function calls
(gdb) disass : Disassemble the code from $eip. An address can be given as an argument.
(gdb) p $eip : print the value of the instruction pointer register.
		The command print can also be used to print variables.
(gdb) set i=42 : Set the variable value to 42.
(gdb) symbole-file kernel : Load the binary which contains the symbol table.
(gdb) info locals : Print data about the local variables.
(gdb) info registers : Print data about the registers
(gdb) backtrace : Print the functions calls to the current point.
\end{verbatim}

\section{Use of both tools}

First of all, to be able to take the advantage of GDB with kaneton, flags must be added to the compilation. In the file environment/env.mk, students are invited to add -g -ggdb to \_CC\_FLAGS\_.

Once it is done, students have to recompile their project and launch it with Qemu and a few more options :
\begin{verbatim}
louis@louis-VirtualBox:/tmp/kaneton$ make clear && make && make build && make install
louis@louis-VirtualBox:/tmp/kaneton$ qemu -s -S -fda kaneton.img
\end{verbatim}

The option -S block the CPU execution. As a consequence, Qemu waits for an order to start the program execution. Moreover, the option -s is an alias for the option -gdb tcp:1234 which creates a socket and wait for a connection from GDB. The communication, once establish between the two program, is made via a pipe (man Qemu, man pipe).

In another terminal, students can launch GDB and attach themself to Qemu which is waiting for a connection :
\begin{verbatim}
louis@louis-VirtualBox:~$ gdb
GNU gdb (Ubuntu/Linaro 7.3-0ubuntu2) 7.3-2011.08
Copyright (C) 2011 Free Software Foundation, Inc.
License GPLv3+: GNU GPL version 3 or later <http://gnu.org/licenses/gpl.html>
This is free software: you are free to change and redistribute it.
There is NO WARRANTY, to the extent permitted by law.  Type "show copying"
and "show warranty" for details.
This GDB was configured as "i686-linux-gnu".
For bug reporting instructions, please see:
<http://bugs.launchpad.net/gdb-linaro/>.
(gdb) target remote localhost:1234
Remote debugging using localhost:1234
0x0000fff0 in ?? ()
(gdb) symbol-file /tmp/kaneton/kaneton/kaneton
Reading symbols from /tmp/kaneton/kaneton/kaneton...done.
\end{verbatim}

The line target remote localhost:1234 establishes the connection with Qemu which is in stand by. More generally, the command target is useful to attach GDB to a program in execution.

When the kernel is installed on the floppy, all symbols are lost. Reloading those symbol is useful to be efficient in debuging. Thus, the command line symbol-file read the ELF File produced by the compilation first step and load all found symbols.

\begin{verbatim}
(gdb) disassemble kaneton
Dump of assembler code for function kaneton:
   0x01000080 <+0>:	add    %al,(%eax)
   0x01000082 <+2>:	add    %al,(%eax)
   0x01000084 <+4>:	add    %al,(%eax)
   0x01000086 <+6>:	add    %al,(%eax)
   0x01000088 <+8>:	add    %al,(%eax)
   0x0100008a <+10>:	add    %al,(%eax)
   0x0100008c <+12>:	add    %al,(%eax)
   0x0100008e <+14>:	add    %al,(%eax)
   0x01000090 <+16>:	add    %al,(%eax)
   ...
\end{verbatim}
Here, the CPU hasn't load the kernel yet due to the -s and -S options. As a consequences, if one get the address of kaneton function, no instruction can yet be interpreted. Thus, by putting a break at the beginning of the function and launching the execution, an interruption is catched and the execution stop at the address 0x01000080.

\begin{verbatim}
(gdb) b *0x01000080
Breakpoint 1 at 0x1000080: file core.c, line 270.
(gdb) c
Continuing.

Breakpoint 1, kaneton (init=0x10b9000) at core.c:270
270	{
\end{verbatim}

Now, the dissamebling can interpret all the symbols and one can get information relatives to the registers.

\begin{verbatim}
(gdb) disass
Dump of assembler code for function kaneton:
=> 0x01000080 <+0>:	push   %ebp
   0x01000081 <+1>:	mov    %esp,%ebp
   0x01000083 <+3>:	sub    $0x38,%esp
   0x01000086 <+6>:	mov    0x8(%ebp),%eax
   0x01000089 <+9>:	mov    %eax,0x105d200
   0x0100008e <+14>:	call   0x1045b90 <module_console_load>
   0x01000093 <+19>:	call   0x1045daf <module_report_load>
   ...
(gdb) info registers
eax            0x10b9000	17534976
ecx            0x205000	2117632
edx            0x1000080	16777344
ebx            0x2c660	181856
esp            0x10c0f78	0x10c0f78
ebp            0x10c0f80	0x10c0f80
esi            0x2cb0e	183054
edi            0x2cb0f	183055
eip            0x1000080	0x1000080 <kaneton>
eflags         0x2	[ ]
cs             0x8	8
ss             0x10	16
ds             0x10	16
es             0x10	16
fs             0x10	16
gs             0x10	16
\end{verbatim}
\chapter{Tests suite}
\label{chapter:Tests suite}

This chapter aims to help students to develop their own tests suite. If some students have another idea of tests suite architecture, they must feel free to implement their own.

The script language used to make this tests suite is Bash. Nevertheless, if students rather use a different one, they may change without a doubt.

\section{First step}
To begin, students have to create the following tree :
\begin{verbatim}
louis@louis-VirtualBox:~$ find testsuite
testsuite
testsuite/utils/
testsuite/tests/
testsuite/tests/k1
testsuite/tests/k2
testsuite/tests/k3
\end{verbatim}

In \textit{utils} directory, students have to copy the core.c of their kaneton. With such an action, they keep the possibility to restore this file as they will have to modify it to launch their tests. Moreover, they have to make a second copy in which they will insert symbols to locate the code insertion in the file.

\begin{verbatim}
louis@louis-VirtualBox:~$ cp kaneton/kaneton/core/core.c testsuite/utils/coreSave.c
louis@louis-VirtualBox:~$ cp kaneton/kaneton/core/core.c testsuite/utils/coreTest.c
\end{verbatim}

Finally, students may insert the symbol '@@@@@' just above the \textit{void kaneton (s\_init* init)} function and '\%\%\%\%\%'  just above the call to \textit{module\_call(test, run)} in 'coreTest.c' file.

\section{Tests suite - 1}
% global structure
\subsection{Global structure}
The goal of this section is to have a tests suite which can execute all tests and produce the results in a text file. The next one will focus on the exploitation of those results.

The idea of this tests suite is to inject some test functions in the core.c and to call them in order to check the kaneton behaviour. Even if a test module exist to allow contributors to check students work, the presented test suite offers the great advantage to the students to masterize their own test environment.

% some test file
\subsection{Test files}
First of all, students must create test files. In this architecture, students only have to code a function (or several if needed, with a main one) in a C file (*.c) without carring about includes.

\begin{verbatim}
louis@louis-VirtualBox:~$ cat testsuite/tests/k1/test1.c

void k1_test1(void)
{
  asm volatile ("int $2\r\n");
}
\end{verbatim}

% Algorithm
\subsection{Algorithm}
If they want, students may create folders in order to be even more explicit in their tests definition. Moreover, as the example proves it, the name of the file can differ from the name of the function. The only constraint is to have no argument for the called function.

Now they have their test files, students can :
\begin{itemize}
\item use find command to create a loop and execute all test files. (-name '*.c')
\item with \textit{grep} and \textit{sed} extract the function name
\item read line by line the coreTest.c file, produce a copy (~/testsuite/coreInjected.c) by replacing the symbole '@@@@@' by the all current test file content and the symbole '\%\%\%\%\%' by a call to the principle function.
\item Replace the kaneton/kaneton/core/core.c file with the coreInjected one.
\end{itemize}

Here follows a fex bash syntax reminder :
\begin{verbatim}
if [ "$?" -eq "0" ];
then
	# Do something
fi

sed 's/.*\(plop\).*/\1/g' # Extract plop from a line
\end{verbatim}

\subsection{Compilation}
% Env modification and compilation
Now, students should be ready to compile their work. Nevertheless, once step must be realised before they do so.

In fact, they could compile but they wouldn't be able to execute their tests suite in nographic mode and being able to recover the output elsewhere. In order to do that, students should activate forward module by modifying \_MODULES\_ line in the environment/env.mk file.

\begin{verbatim}
_MODULES_ := console forward
\end{verbatim}

The forward module is used to redirect the SDL output over a serial port. Then, students will be able redirect that flow to the standard unix output by adding options to qemu command line.

\textbf{N.B. : Students may activate it automatically using sed program.}

Now students may use cd to go to the kaneton directory and compile it with the usual compilation line.

\subsection{Qemu}
To launch their kaneton and being able to redirect the flow to the standard output or in a file, students have to add options to their call to Qemu.

In addition to  the option
\begin{verbatim} -fda kaneton.img \end{verbatim}
, students need to add the option :
\begin{verbatim} -nographic \end{verbatim}
The option nographic include a redirection of the serial port to the standard unix output (-serial stdio) and disable all graphic output (such as SDL).

As a conclusion, students may want to use the following command line :
\begin{verbatim}
qemu -nographic -fda kaneton.img >> /desired/path/outputFile.txt
\end{verbatim}


Those are the objective of the next section.

\section{Tests suite - 2}
In the previous section, students where guided to have a program which can execute kaneton with the desired test function and redirect the output to a file. In that section, students won't be as guided as before. Here, they will just have a guideline to compelete their tests suite.


Everybody will notice at this point that two things are missing. Firstly, kaneton finishes with an infinit loop. Thus, it will be mandatory to managed that end. Two cases are possible : the kernel has finished its execution and doesn't launch any server or the test creates a timetimout. To manage that difficulty, students may think about what is printed at the end of the kaneton execution and how can they simulate a fork.

Secondly, the produced output have to be exloited to check if the test can be validated or not. To reach that goal, students are invited to create a fileTest.res, next to their fileTest.c, and to compare the expected results (written in fileTest.res) with the actual results (/desired/path/outputFile.txt). That comparison is heart of the test validation.

\chapter{Conclusion}
Students are supposed to find everything they need in those assignments to feel comfortable to debug their kaneton.

However, as it as been said in introduction, comments are welcome to improve that suject which is a few days old.

Good luck with the end of k1 ;).

%
% ---------- header -----------------------------------------------------------
%
% project       kaneton
%
% license       kaneton
%
% file          /home/mycure/kaneton/view/book/core/licenses.tex
%
% created       julien quintard   [tue may 12 18:35:31 2009]
% updated       julien quintard   [tue may 12 18:35:33 2009]
%

%
% ---------- licenses ---------------------------------------------------------
%

\chapter{Licenses}
\label{chapter:licenses}

In this chapter we will details the kaneton-related licenses.

\newpage

%
% ---------- text -------------------------------------------------------------
%

The kaneton project might be considered as an open-project since source
code is provided.

Nevertheless this is not the case as this project is used as material for
operating system courses.

Therefore, people implementing the kaneton microkernel should not
make their source code available. To avoid problems, especially students
cheating, kaneton people decided to use a kaneton-specific license
forbidding source code distribution.

The \name{kaneton license} is based on a more generic license, the
\name{pedagogical licence}.

The next sections will contain these licenses' descriptions.

%
% pedagogical license
%

\section{Pedagogical License}

\textbf{TERMS AND CONDITIONS FOR COPYING, DISTRIBUTION AND MODIFICATION}

\begin{enumerate}
  \item
    This License applies to any program or other work which contains a
    notice placed by the copyright holder saying it may be distributed
    under the terms of this License. The "Program", below, refers
    to any such program or work, and a "work based on the Program"
    means either the Program or any derivative work under copyright law:
    that is to say, a work containing the Program or a portion of it,
    either verbatim or with modifications and/or translated into another
    language. (Hereinafter, translation is included without limitation in
    the term "modification".) Each licensee is addressed as "you".
  \item
    You must not copy or distribute copies of the Program's source
    code, object code or executable form without explicit authorization
    from the maintainers.

    If you have this authorization, you must conspicuously and
    appropriately publish on each copy an appropriate copyright notice
    and disclaimer of warranty; keep intact all the notices that refer
    to this License and to the absence of any warranty; and give any
    other recipients of the Program a copy of this License along with
    the Program.
  \item
    You may modify your copy or copies of the Program or any portion of
    it, thus forming a work based on the Program, provided that you
    meet all of these conditions :

    \begin{enumerate}
      \item
	You must not publish your work without explicit authorization
	from the maintainers.
      \item
	You must send to the maintainers any work that in whole or in
	part contains or is derived from the Program or any part thereof.
      \item
	You must cause any work that you send, that in whole or in
	part contains or is derived from the Program or any part thereof,
	to be licensed as a whole under the terms of this License.
      \item
	You must cause the modified files to carry prominent notices
	stating that you changed the files and the date of any change.
      \item
	If the modified program normally reads commands interactively
	when run, you must cause it, when started running for such
	interactive use in the most ordinary way, to print or display an
	announcement including an appropriate copyright notice and a
	notice that there is no warranty (or else, saying that you
	provide a warranty) and telling the user how to view a copy
	of this License. (Exception: if the Program itself is interactive
	but does not normally print such an announcement, your work based
	on the Program is not required to print an announcement.)
    \end{enumerate}
  \item
    Access to the Program's source is granted if either:

    \begin{enumerate}
      \item
	You want to make the Program evolve
      \item
	You have pedagogical goals
    \end{enumerate}

    These requirements apply to the modified work as a whole. If
    identifiable sections of that work are not derived from the Program,
    and can be reasonably considered independent and separate works in
    themselves, then this License, and its terms, do not apply to those
    sections when you distribute them as separate works. But when you
    distribute the same sections as part of a whole which is a work based
    on the Program, the distribution of the whole must be on the terms of
    this License, whose permissions for other licensees extend to the
    entire whole, and thus to each and every part regardless of who wrote
    it.

    Thus, it is not the intent of this section to claim rights or contest
    your rights to work written entirely by you; rather, the intent is to
    exercise the right to control the distribution of derivative or
    collective works based on the Program.

    In addition, mere aggregation of another work not based on the Program
    with the Program (or with a work based on the Program) on a volume of
    a storage or distribution medium does not bring the other work under
    the scope of this License.
  \item
    You may not copy, modify, sublicense, or distribute the Program
    except as expressly provided under this License. Any attempt
    otherwise to copy, modify, sublicense or distribute the Program is
    void, and will automatically terminate your rights under this
    License. However, parties who have received copies, or rights,
    from you under this License will not have their licenses
    terminated so long as such parties remain in full compliance.
  \item
    You are not required to accept this License, since you have not
    signed it. However, nothing else grants you permission to modify
    or distribute the Program or its derivative works. These actions
    are prohibited by law if you do not accept this
    License. Therefore, by modifying or distributing the Program (or
    any work based on the Program), you indicate your acceptance of
    this License to do so, and all its terms and conditions for
    copying, distributing or modifying the Program or works based on
    it.
  \item
    Each time you redistribute the Program (or any work based on the
    Program), the recipient automatically receives a license from the
    original licensor to copy, distribute or modify the Program subject
    to these terms and conditions. You may not impose any further
    restrictions on the recipients' exercise of the rights granted
    herein. You are not responsible for enforcing compliance by third
    parties to this License.
  \item
    If the distribution and/or use of the Program is restricted in
    certain countries either by patents or by copyrighted interfaces,
    the original copyright holder who places the Program under this
    License may add an explicit geographical distribution limitation
    excluding those countries, so that distribution is permitted only
    in or among countries not thus excluded. In such case, this License
    incorporates the limitation as if written in the body of this
    License.
  \item
    We may publish revised and/or new versions of this License from
    time to time. It may evolve considering new contributors
    needs. Contact us if you have any request. Such new versions will
    be similar in spirit to the present version, but may differ in
    detail to address new problems or concerns.
  \item
    If you wish to incorporate parts of the Program into other free
    programs whose distribution conditions are different, write to the
    author to ask for permission, we sometimes make exceptions for this.
  \item
    THERE IS NO WARRANTY FOR THE PROGRAM, TO THE EXTENT PERMITTED BY
    APPLICABLE LAW. EXCEPT WHEN OTHERWISE STATED IN WRITING THE
    COPYRIGHT HOLDERS AND/OR OTHER PARTIES PROVIDE THE PROGRAM "AS IS"
    WITHOUT WARRANTY OF ANY KIND, EITHER EXPRESSED OR IMPLIED,
    INCLUDING, BUT NOT LIMITED TO, THE IMPLIED WARRANTIES OF
    MERCHANTABILITY AND FITNESS FOR A PARTICULAR PURPOSE. THE ENTIRE
    RISK AS TO THE QUALITY AND PERFORMANCE OF THE PROGRAM IS WITH
    YOU. SHOULD THE PROGRAM PROVE DEFECTIVE, YOU ASSUME THE COST OF
    ALL NECESSARY SERVICING, REPAIR OR CORRECTION.
  \item
    IN NO EVENT UNLESS REQUIRED BY APPLICABLE LAW OR AGREED TO IN
    WRITING WILL ANY COPYRIGHT HOLDER, OR ANY OTHER PARTY WHO MAY
    MODIFY AND/OR REDISTRIBUTE THE PROGRAM AS PERMITTED ABOVE, BE
    LIABLE TO YOU FOR DAMAGES, INCLUDING ANY GENERAL, SPECIAL,
    INCIDENTAL OR CONSEQUENTIAL DAMAGES ARISING OUT OF THE USE OR
    INABILITY TO USE THE PROGRAM (INCLUDING BUT NOT LIMITED TO LOSS OF
    DATA OR DATA BEING RENDERED INACCURATE OR LOSSES SUSTAINED BY YOU
    OR THIRD PARTIES OR A FAILURE OF THE PROGRAM TO OPERATE WITH ANY
    OTHER PROGRAMS), EVEN IF SUCH HOLDER OR OTHER PARTY HAS BEEN
    ADVISED OF THE POSSIBILITY OF SUCH DAMAGES.
\end{enumerate}

%
% kaneton license
%

\section{kaneton License}

\textbf{TERMS AND CONDITIONS FOR COPYING, DISTRIBUTION AND MODIFICATION}

\begin{enumerate}
  \item
    This licence is nothing more than a link to the pedagogical licence.
  \item
    Any program under the kaneton licence is in fact under the terms and
    conditions of the pedagogical licence.
\end{enumerate}


\end{document}
